\section{引言：一个激进的实验}

2025 年 8 月，OpenAI 内部一个小团队启动了一项激进的实验：从一个空的 Git 仓库开始，构建一个完整的软件产品，且\textbf{不写一行人工代码}。五个月后，这个仓库包含了约一百万行代码，涵盖应用逻辑、基础设施、工具链、文档以及内部开发者工具。在此期间，一个仅三人的工程团队通过 Codex 提交并合并了约 1,500 个 Pull Request，平均每人每天 3.5 个 PR。该产品已在 OpenAI 内部拥有数百名用户，包括每日活跃的重度用户。

\begin{importantbox}{核心理念：Harness Engineering}
Harness Engineering（驾驭工程）是文章提出的核心概念。当 AI Agent 成为主要的代码产出者时，软件工程师的角色从``写代码''转变为``设计环境、表达意图、构建反馈循环''。人类驾驭（harness）Agent，而非亲自执行编码任务。这个名字来源于马具（harness）的隐喻：你不替马跑，而是为它装上缰绳和方向盘。
\end{importantbox}

这篇文章是 OpenAI 内部团队的一线实战报告，记录了在完全 Agent 驱动的开发模式下遇到的问题、积累的经验、以及逐步形成的方法论。它回答了一个关键问题：\textbf{当工程师不再写代码时，他们应该做什么？}

\begin{knowledgebox}{背景：Codex 与 Agent-First 开发}
Codex 是 OpenAI 推出的 coding agent 产品，能够理解自然语言指令并生成代码、运行测试、提交 PR。``Agent-First''开发范式指的是将 AI Agent 作为主要的代码生产者，人类工程师转而负责系统设计、质量保障和方向把控。这与传统的``AI 辅助编程''（如 Copilot 自动补全）有本质区别：后者是人写代码、AI 辅助，前者是 AI 写代码、人类引导。
\end{knowledgebox}

\subsection{实验数据概览}

\begin{table}[H]
\centering
\begin{tabular}{ll}
\toprule
\textbf{指标} & \textbf{数据} \\
\midrule
项目起始时间 & 2025 年 8 月底 \\
代码规模 & 约 100 万行 \\
PR 总数 & 约 1,500 个 \\
初始团队规模 & 3 名工程师 \\
后续团队规模 & 7 名工程师 \\
人均日 PR 数 & 3.5 个 \\
人工代码行数 & 0 行 \\
估算提速倍数 & 约 10 倍 \\
单次 Codex 最长运行时间 & 6 小时以上 \\
\bottomrule
\end{tabular}
\caption{OpenAI Harness Engineering 实验关键数据}
\end{table}

\subsection{本章小结}

OpenAI 团队通过一个极端约束（零人工代码）验证了 Agent-First 开发范式的可行性。关键发现是：Agent 产出代码的速度远超人类，但\textbf{真正的瓶颈不在于代码生成，而在于人类的注意力和判断力}。整篇文章围绕如何最大化利用这一稀缺资源展开。

\newpage

